% Capítulo — TCC 1
\chapter{Metodologia}
\label{chap:metodologia}

% Abertura do capítulo: um parágrafo dizendo o que ele apresenta.

\section{Design Science Research com estudo de caso}
\label{sec:metodologia-design-science-research-com-estudo-de-caso}
% Hevner et al. (2004); Peffers et al. (2007); ciclos de construção e avaliação; estudo de caso com bases do GovHub.

\section{Spec-Driven Development}
\label{sec:metodologia-spec-driven-development}
% Ciclo: spec com critérios de aceite, plano, implementação por agente de IA, validação, ADR. Modelo em docs/specs/_template.md.

\section{Ciclos de iteração}
\label{sec:metodologia-ciclos-de-iteracao}
% PoC IBGE, costuras A/B/C, avaliação; evidência no histórico do repositório.

\section{Ground truth}
\label{sec:metodologia-ground-truth}
% Seleção de 20 a 40 pares, três níveis de dificuldade, rótulo por categoria de atrito, pares negativos, congelamento antes de ajustar o agente.

\section{Protocolo de avaliação}
\label{sec:metodologia-protocolo-de-avaliacao}
% Tabela QP x experimento x métricas x baseline; 3 a 5 execuções por par com LLM; registro das specs para a QP5.

\section{Ameaças à validade}
\label{sec:metodologia-ameacas-a-validade}
% Interna, externa, de constructo e de conclusão (Wohlin et al., 2012).

\section{Cronograma}
\label{sec:metodologia-cronograma}
% Fases até a defesa do TCC 2 (agosto de 2027); portões: ground truth congelado, resultados congelados.
